% !TeX encoding = UTF-8
% 字体映射。此文件由主文件在 fontspec、xeCJK 加载后读入。
% 优先使用本机已安装的原文字体；缺少时采用 fonts/ 中的开源替代字体。
% 所有字体名、样式和图片以原文为依据；替代字体并非原字体的逐字形副本。

% 西文主字体与 Times New Roman 对应的文字运行。
\IfFontExistsTF{Times New Roman}{%
  \setmainfont{Times New Roman}
  \newfontfamily\WordLatinTimes{Times New Roman}
}{%
  \setmainfont[
    Path=fonts/,
    BoldFont=texgyretermes-bold.otf,
    ItalicFont=texgyretermes-italic.otf,
    BoldItalicFont=texgyretermes-bolditalic.otf
  ]{texgyretermes-regular.otf}
  \newfontfamily\WordLatinTimes[
    Path=fonts/,
    BoldFont=texgyretermes-bold.otf,
    ItalicFont=texgyretermes-italic.otf,
    BoldItalicFont=texgyretermes-bolditalic.otf
  ]{texgyretermes-regular.otf}
}

% 宋体：同时映射西文运行及中文运行。
\IfFontExistsTF{SimSun}{%
  \setCJKmainfont[AutoFakeBold=1.5,AutoFakeSlant=.2]{SimSun}
  \setCJKfamilyfont{WordSong}[AutoFakeBold=1.5,AutoFakeSlant=.2]{SimSun}
  \newfontfamily\WordLatinSong[AutoFakeBold=1.5,AutoFakeSlant=.2]{SimSun}
}{%
  \setCJKmainfont[
    Path=fonts/,Script=Default,
    BoldFont=FandolSong-Bold.otf,
    ItalicFont=FandolKai-Regular.otf
  ]{FandolSong-Regular.otf}
  \setCJKfamilyfont{WordSong}[
    Path=fonts/,Script=Default,
    BoldFont=FandolSong-Bold.otf,
    ItalicFont=FandolKai-Regular.otf
  ]{FandolSong-Regular.otf}
  \newfontfamily\WordLatinSong[
    Path=fonts/,Script=Default,
    BoldFont=FandolSong-Bold.otf,
    ItalicFont=FandolKai-Regular.otf
  ]{FandolSong-Regular.otf}
}

% 黑体。
\IfFontExistsTF{SimHei}{%
  \setCJKfamilyfont{WordHei}[AutoFakeBold=1.5,AutoFakeSlant=.2]{SimHei}
  \newfontfamily\WordLatinHei[AutoFakeBold=1.5,AutoFakeSlant=.2]{SimHei}
}{%
  \setCJKfamilyfont{WordHei}[
    Path=fonts/,Script=Default,
    BoldFont=FandolHei-Bold.otf,AutoFakeSlant=.2
  ]{FandolHei-Regular.otf}
  \newfontfamily\WordLatinHei[
    Path=fonts/,Script=Default,
    BoldFont=FandolHei-Bold.otf,AutoFakeSlant=.2
  ]{FandolHei-Regular.otf}
}

% 楷体：兼容 Windows 新旧名称。
\IfFontExistsTF{KaiTi}{%
  \setCJKfamilyfont{WordKai}[AutoFakeBold=1.5,AutoFakeSlant=.2]{KaiTi}
  \newfontfamily\WordLatinKai[AutoFakeBold=1.5,AutoFakeSlant=.2]{KaiTi}
}{%
  \IfFontExistsTF{KaiTi_GB2312}{%
    \setCJKfamilyfont{WordKai}[AutoFakeBold=1.5,AutoFakeSlant=.2]{KaiTi_GB2312}
    \newfontfamily\WordLatinKai[AutoFakeBold=1.5,AutoFakeSlant=.2]{KaiTi_GB2312}
  }{%
    \setCJKfamilyfont{WordKai}[
      Path=fonts/,Script=Default,AutoFakeBold=1.5,AutoFakeSlant=.2
    ]{FandolKai-Regular.otf}
    \newfontfamily\WordLatinKai[
      Path=fonts/,Script=Default,AutoFakeBold=1.5,AutoFakeSlant=.2
    ]{FandolKai-Regular.otf}
  }
}

% 仿宋。
\IfFontExistsTF{FangSong}{%
  \setCJKfamilyfont{WordFang}[AutoFakeBold=1.5,AutoFakeSlant=.2]{FangSong}
  \newfontfamily\WordLatinFang[AutoFakeBold=1.5,AutoFakeSlant=.2]{FangSong}
}{%
  \setCJKfamilyfont{WordFang}[
    Path=fonts/,Script=Default,AutoFakeBold=1.5,AutoFakeSlant=.2
  ]{FandolFang-Regular.otf}
  \newfontfamily\WordLatinFang[
    Path=fonts/,Script=Default,AutoFakeBold=1.5,AutoFakeSlant=.2
  ]{FandolFang-Regular.otf}
}

% 华文中宋。Fandol 没有同字形字体，缺少原字体时用宋体风格替代。
\IfFontExistsTF{STZhongsong}{%
  \setCJKfamilyfont{WordZhong}[AutoFakeBold=1.5,AutoFakeSlant=.2]{STZhongsong}
  \newfontfamily\WordLatinZhong[AutoFakeBold=1.5,AutoFakeSlant=.2]{STZhongsong}
}{%
  \setCJKfamilyfont{WordZhong}[
    Path=fonts/,Script=Default,AutoFakeBold=1.5,AutoFakeSlant=.2
  ]{FandolSong-Regular.otf}
  \newfontfamily\WordLatinZhong[
    Path=fonts/,Script=Default,AutoFakeBold=1.5,AutoFakeSlant=.2
  ]{FandolSong-Regular.otf}
}

% 少量继承自 Word 默认样式的西文字体。
\IfFontExistsTF{Cambria}{%
  \newfontfamily\WordLatinCambria{Cambria}
}{%
  \let\WordLatinCambria\WordLatinTimes
}
\IfFontExistsTF{Calibri}{%
  \newfontfamily\WordLatinCalibri{Calibri}
}{%
  \let\WordLatinCalibri\WordLatinTimes
}
